% Results and Discussion: VERACISM

\section{RESULTS AND DISCUSSION}

This section introduces and explains the results of the five testing categories
that have been performed to assess the VERACISM platform: (1) the functional testing,
(2) the integration test, (3) the performance testing, (4) the security testing, and (5) the user acceptance testing.

The appendices show the raw data for each category completely. All the
assertion-level tables are detailed in the integration tests and functional tests
Appendices 0.9 and 0.10. The security Appendix 0.11 includes security findings. Performance configuration is part of the performance Appendix 0.12. And the functional through security remediation PDF appendices contain full original test report PDFs.

% ----------------------------------------------------------
\subsection{Functional and Integration Testing}
% ----------------------------------------------------------

The hybrid method outlined in the Project Assessment chapter – automated black box API testing using Newman and manual role based functional verification – was used for functional and integration testing. The two streams shared the same scenario set, with those scenarios scripted as automation steps and others performed manually (for their dependency on business logic and/or user interface considerations).

\subsubsection{The automated test summary}

The automated Integration Test Suite consisted of five test groups that interact with the external VERACISM API v1: scan upload, report retrieval by CID, public verification, controlled download and storage status check. Twenty two (22) test cases with a total of 93 individual assertions were run against the staging environment in \url{https://sb-api.ismanila.org/Veracism}. Four (4) user role groups were covered in 19 automated tests defined by 89 assertions in the Functional Test Suite. Everyone received a pass mark of 100\% for both suites and there were no failures or errors.

Table~\ref{tab:combined-test-summary} shows the combined statistics.

\begin{table}[H]
  \vspace{2ex}
  \centering
  \footnotesize
  {%
  \setlength{\tabcolsep}{6pt}
  \renewcommand{\arraystretch}{1.3}
  \begin{tabular}{|>{\raggedright\arraybackslash}p{0.30\textwidth}|
                  >{\centering\arraybackslash}p{0.13\textwidth}|
                  >{\centering\arraybackslash}p{0.13\textwidth}|
                  >{\centering\arraybackslash}p{0.10\textwidth}|
                  >{\centering\arraybackslash}p{0.10\textwidth}|
                  >{\centering\arraybackslash}p{0.13\textwidth}|}
    \hline
    \textbf{Suite} & \textbf{Test Cases} & \textbf{Assertions} & \textbf{Passed} & \textbf{Failed} & \textbf{Pass Rate} \\
    \hline\hline
    Integration Test Suite & 22 & 93 & 93 & 0 & 100.0\% \\
    \hline
    Functional Test Suite & 19 & 89 & 89 & 0 & 100.0\% \\
    \hline
    \textbf{Combined Total} & \textbf{41} & \textbf{182} & \textbf{182} & \textbf{0} & \textbf{100.0\%} \\
    \hline
  \end{tabular}
  }%
  \vspace{2ex}
  \captionsetup{type=table}
  \caption{Combined automated test results: Integration and Functional suites (30--31 March 2026)}
  \label{tab:combined-test-summary}
\end{table}

\subsubsection{Integration suite by endpoint group}

Table~\ref{tab:int-group-breakdown} represents the 93~integration assertions run in the five endpoint groups, along with the overall run time of the Newman runner. Most of the elapsed time is spent in the Scan Upload group, representing 83.7\% of the total run time of 83.7 seconds; each of these ``scan upload'' calls triggers the whole document pipe: SHA-256 computation, ClamAV scanning, YARA scanning, Lighthouse encryption and IPFS anchoring.

\begin{table}[H]
  \vspace{2ex}
  \centering
  \footnotesize
  {%
  \setlength{\tabcolsep}{6pt}
  \renewcommand{\arraystretch}{1.3}
  \begin{tabular}{|>{\raggedright\arraybackslash}p{0.29\textwidth}|
                  >{\centering\arraybackslash}p{0.09\textwidth}|
                  >{\centering\arraybackslash}p{0.11\textwidth}|
                  >{\centering\arraybackslash}p{0.07\textwidth}|
                  >{\centering\arraybackslash}p{0.07\textwidth}|
                  >{\centering\arraybackslash}p{0.11\textwidth}|
                  >{\centering\arraybackslash}p{0.10\textwidth}|}
    \hline
    \textbf{Endpoint Group} & \textbf{TCs} & \textbf{Assertions} & \textbf{Pass} & \textbf{Fail} & \textbf{Time (s)} & \textbf{\% of Run} \\
    \hline\hline
    Scan Upload & 8 & 37 & 37 & 0 & 77.969 & 93.1\% \\
    \hline
    Get Report by CID & 4 & 17 & 17 & 0 & 0.307 & 0.4\% \\
    \hline
    Verify (Public) & 4 & 18 & 18 & 0 & 0.230 & 0.3\% \\
    \hline
    Download & 3 & 12 & 12 & 0 & 6.014 & 7.2\% \\
    \hline
    Storage Status & 3 & 11 & 11 & 0 & 0.163 & 0.2\% \\
    \hline
    \textbf{TOTAL} & \textbf{22} & \textbf{93} & \textbf{93} & \textbf{0} & \textbf{83.728} & \textbf{100\%} \\
    \hline
  \end{tabular}
  }%
  \vspace{2ex}
  \captionsetup{type=table}
  \caption{Grouped by endpoints, the following is the result of running the integration test suite and providing full assertion details in Appendix~\ref{appendix:integration-tests})}
  \label{tab:int-group-breakdown}
\end{table}

\noindent The breakdown brings up an important aspect, the page load time -- every time the metadata is retrieved from the file (which also includes a lookup for its CID with 0.147 seconds -- still in interactive times, despite sending the message cross-network), anyone is able to verify its currency (0.066 seconds -- again, within interactive expectation), and one can check if it is stored or not (0.061 seconds -- still within interactive expectations). The upload pipeline is inherently known to be capped by the round trip to the Lighthouse encryption proxy and the Filecoin network - in this case of 15-21 seconds per file, since the round trip can't be avoided. This is around the normal latency of file anchoring operations with Lighthouse API.



\subsection{Functional suite by automated group}
Table~\ref{tab:func-auto-breakdown} lists the 89 functional assertions for four of the Newman runner groups: F1, IT Office API; F2, Registrar Upload; F3, Public Verification; F4, End-to-End Journey.
The total time spent on automated files was \textbf{116.641~seconds}.

\begin{table}[H]
  \vspace{2ex}
  \centering
  \footnotesize
  {%
  \setlength{\tabcolsep}{6pt}
  \renewcommand{\arraystretch}{1.3}
  \begin{tabular}{|>{\raggedright\arraybackslash}p{0.29\textwidth}|
                  >{\centering\arraybackslash}p{0.09\textwidth}|
                  >{\centering\arraybackslash}p{0.11\textwidth}|
                  >{\centering\arraybackslash}p{0.07\textwidth}|
                  >{\centering\arraybackslash}p{0.07\textwidth}|
                  >{\centering\arraybackslash}p{0.11\textwidth}|
                  >{\centering\arraybackslash}p{0.10\textwidth}|}
    \hline
    \textbf{Feature Group} & \textbf{Suites} & \textbf{Assertions} & \textbf{Pass} & \textbf{Fail} & \textbf{Time (s)} & \textbf{\% of Run} \\
    \hline\hline
    F1: IT Office API & 5 & 27 & 27 & 0 & 54.231 & 46.5\% \\
    \hline
    F2: Registrar Upload & 4 & 20 & 20 & 0 & 39.508 & 33.9\% \\
    \hline
    F3: Public Verification & 7 & 28 & 28 & 0 & 21.047 & 18.0\% \\
    \hline
    F4: End-to-End Journey & 3 & 14 & 14 & 0 & 1.855 & 1.6\% \\
    \hline
    \textbf{TOTAL} & \textbf{19} & \textbf{89} & \textbf{89} & \textbf{0} & \textbf{116.641} & \textbf{100\%} \\
    \hline
  \end{tabular}
  }%
  \vspace{2ex}
  \captionsetup{type=table}
  \caption{Functional test suite automated results by feature group (31 March 2026)}
  \label{tab:func-auto-breakdown}
\end{table}

\subsubsection{Manual and code review results – per role}

In total, there are 30 test cases (automated API, manual code review and observational inspection) across all five user roles in the system specification.
The pass, partial and fail split for each role group is shown in Table~\ref{tab:func-manual-by-role} 

\begin{table}[H]
  \vspace{2ex}
  \centering
  \footnotesize
  {%
  \setlength{\tabcolsep}{6pt}
  \renewcommand{\arraystretch}{1.3}
  \begin{tabular}{|>{\raggedright\arraybackslash}p{0.28\textwidth}|
                  >{\centering\arraybackslash}p{0.08\textwidth}|
                  >{\centering\arraybackslash}p{0.09\textwidth}|
                  >{\centering\arraybackslash}p{0.09\textwidth}|
                  >{\centering\arraybackslash}p{0.09\textwidth}|
                  >{\centering\arraybackslash}p{0.09\textwidth}|
                  >{\centering\arraybackslash}p{0.09\textwidth}|}
    \hline
    \textbf{Role Group} & \textbf{Total} & \textbf{PASSED} & \textbf{PARTIAL} & \textbf{FAILED} & \textbf{N/T} & \textbf{Pass \%} \\
    \hline\hline
    System Administrator (TC-ADM) & 5 & 5 & 0 & 0 & 0 & 100.0\% \\
    \hline
    IT Office Staff (TC-ITO) & 6 & 4 & 1 & 1 & 0 & 66.7\% \\
    \hline
    Registrar / Admin Office (TC-REG) & 7 & 7 & 0 & 0 & 0 & 100.0\% \\
    \hline
    Student (TC-STU) & 5 & 5 & 0 & 0 & 0 & 100.0\% \\
    \hline
    External Auditor / Educator (TC-EXT) & 4 & 4 & 0 & 0 & 0 & 100.0\% \\
    \hline
    Cross-Role End-to-End (TC-E2E) & 3 & 2 & 1 & 0 & 0 & 66.7\% \\
    \hline
    \textbf{OVERALL} & \textbf{30} & \textbf{27} & \textbf{2} & \textbf{1} & \textbf{0} & \textbf{90.0\%} \\
    \hline
  \end{tabular}
  }%
  \vspace{2ex}
  \captionsetup{type=table}
  \caption{Functional test results by user role (N/T = not tested in this cycle; full detail in Appendix B, Tables~\ref{tab:func-adm} through~\ref{tab:func-e2e})}
  \label{tab:func-manual-by-role}
\end{table}

\noindent A \textbf{100\% pass} with no failures and partial results was seen by three out of four groups that play the most critical role in the value proposition of the system (System Administrator, Registrar/Admin Office, External Auditor). The single failure (TC-ITO-005) and one partial result (TC-ITO-006) produced the lowest rates for the IT Office group, 66.7\%.

\subsubsection{Open defects summary}

The Manual and code review phase was able to surface three defects.
Details regarding each defect, its severity and disposition are summarized in Table~\ref{tab:func-defects}.

\begin{table}[H]
  \vspace{2ex}
  \centering
  \footnotesize
  {%
  \setlength{\tabcolsep}{5pt}
  \renewcommand{\arraystretch}{1.3}
  \begin{tabular}{|>{\centering\arraybackslash}p{0.12\textwidth}|
                  >{\centering\arraybackslash}p{0.08\textwidth}|
                  >{\raggedright\arraybackslash}p{0.18\textwidth}|
                  >{\raggedright\arraybackslash}p{0.34\textwidth}|
                  >{\centering\arraybackslash}p{0.12\textwidth}|}
    \hline
    \textbf{TC ID} & \textbf{Severity} & \textbf{Title} & \textbf{Description} & \textbf{Disposition} \\
    \hline\hline
    TC-ITO-005 & Medium & No Lighthouse retry policy & No Polly retry or circuit breaker on \texttt{LighthouseStorageService}; transient unavailability causes immediate failure with no backoff & Future Work \\
    \hline
    TC-ITO-006 & Low & Scoped audit log UI incomplete & \texttt{AuditEvent} table has no C\# model, repository, or controller; tenant scoped audit UI cannot be rendered & Future Work \\
    \hline
    TC-E2E-003 & Medium & VerificationLog not written & \texttt{VerificationLog} table exists in schema but \texttt{VerifyController} does not persist a log entry on each \texttt{GET /verify} call & Future Work \\
    \hline
  \end{tabular}
  }%
  \vspace{2ex}
  \captionsetup{type=table}
  \caption{Open defects after functional testing: All of them are to be prioritized for upcoming development sprint}
  \label{tab:func-defects}
\end{table}

\subsubsection{Key functional findings}

The following system behaviours were confirmed by the combined testing programme:

\begin{enumerate}
  \item \textbf{Core pipeline is fully end to end functional.} For TC-E2E-001, the threat scan is successful after completely uploading a file with verified attempt from unauthenticated public verification, \texttt{sha256Match: true} and \texttt{inLighthouse: true}.

  \item \textbf{Multi-engine threat detection gates work properly.} A JPG with a webshell signature is detected through the usage of the YARA engine embedded in the rule: \texttt{image\_webshell\_signature} and cannot be uploaded to Lighthouse, whilst ClamAV says it is clean, as shown above, this multi engine approach proves useful.

  \item \textbf{The SHA-256 integrity is maintained end to end.}  The hash
    computed at upload (TC-ITO-001) is retrievable via the report endpoint
    (TC-ITO-003) and confirmed by the public verification endpoint
    (TC-E2E-001), covering the full document lifecycle.

  \item \textbf{An enumeration is not possible because tenant data isolation.}  TC-11 established that a CID of a different leaseholder will return a 404 response if an approved attacker tries to access a cross tenant CID he can conclude that documents does not exist, not that they are prohibited.

  \item \textbf{The authentication and authorisation boundaries are proper.}  All protected endpoints returned bad HTTP 401/403 responses for requests with invalid or missing tokens. The public verification endpoint returned HTTP 200 without any credentials, TC-EXT-001 and TC-15.

  \item \textbf{Document immutability is enforced at the API level.}  No
    PUT or DELETE endpoints exist for report records (TC-REG-006); the API
    is append only by design, matching the tamper evident guarantee.

  \item \textbf{Tampered-document detection is reliable.}  A fabricated CID
    (TC-E2E-002) returned HTTP~404 with a descriptive message and no
    false positive verification result.
\end{enumerate}

% ----------------------------------------------------------
\subsection{Performance Testing}
% ----------------------------------------------------------

Performance testing was conducted using Postman's Performance Testing feature
targeting the public CID verification endpoint
(\texttt{GET /api/v1/verify?cid=\{cid\}}), which is the endpoint most
exposed to external concurrent traffic in the production workload.  Two
rounds were executed, differing only in whether the Nginx reverse proxy
rate limiting policy was active.

\subsubsection{Test configuration and scope}

Both rounds used 100 virtual users (VUs) sustained for a 5-minute (300-second)
window, meaning the test engine attempted to drive \textbf{up to 30,000
requests} during each run at sustained concurrency.
Table~\ref{tab:perf-config-detail} summarises the configuration of both rounds.

\begin{table}[H]
  \vspace{2ex}
  \centering
  \footnotesize
  {%
  \setlength{\tabcolsep}{6pt}
  \renewcommand{\arraystretch}{1.3}
  \begin{tabular}{|>{\raggedright\arraybackslash}p{0.31\textwidth}|
                  >{\centering\arraybackslash}p{0.26\textwidth}|
                  >{\centering\arraybackslash}p{0.26\textwidth}|}
    \hline
    \textbf{Parameter} & \textbf{Version 1.0 (2026-03-28)} & \textbf{Version 1.1 (2026-03-30)} \\
    \hline\hline
    Target endpoint & \multicolumn{2}{c|}{\texttt{GET /api/v1/verify?cid=\{cid\}}} \\
    \hline
    Virtual users (VUs) & 100 & 100 \\
    \hline
    Test duration & 5 minutes & 5 minutes \\
    \hline
    Infrastructure rate limiting & \textbf{Enabled} (Nginx default) & \textbf{Disabled} (rate-limit rules removed) \\
    \hline
    Error source identified & HTTP~429 Too Many Requests from Nginx proxy & No infrastructure throttle \\
    \hline
    Application layer result & Not reachable under sustained load & Successfully handled concurrent load \\
    \hline
  \end{tabular}
  }%
  \vspace{2ex}
  \captionsetup{type=table}
  \caption{Performance test configuration and run conditions}
  \label{tab:perf-config-detail}
\end{table}

\subsubsection{Key performance findings}

Table~\ref{tab:perf-findings} presents the principal observations and their
implications for production deployment.

\begin{table}[H]
  \vspace{2ex}
  \centering
  \footnotesize
  {%
  \setlength{\tabcolsep}{5pt}
  \renewcommand{\arraystretch}{1.3}
  \begin{tabular}{|>{\raggedright\arraybackslash}p{0.25\textwidth}|
                  >{\raggedright\arraybackslash}p{0.26\textwidth}|
                  >{\raggedright\arraybackslash}p{0.26\textwidth}|
                  >{\raggedright\arraybackslash}p{0.13\textwidth}|}
    \hline
    \textbf{Metric / Observation} & \textbf{Version 1.0 Result} & \textbf{Version 1.1 Result} & \textbf{Root Cause} \\
    \hline\hline
    Error rate & \textbf{Significantly elevated}; dominant error was HTTP~429 (rate-limit exceeded) returned by Nginx before requests reached the backend & \textbf{Significantly reduced}; errors largely eliminated once Nginx throttle removed & Nginx rate-limit threshold set below 100-VU load level \\
    \hline
    Average response time & Higher; most latency was queuing delay at the proxy level, not application processing time & Lower; requests reached the .NET~8 backend directly with minimal queuing overhead & Infrastructure configuration, not application code \\
    \hline
    Throughput (req/s) & Constrained by HTTP~429 rejection at proxy; effective throughput well below 100-VU potential & Higher and stable; application layer did not become a bottleneck during the 5-minute run & Rate-limit rules \\
    \hline
    Application layer & Healthy; .NET~8 backend and SQL~Server correctly processed all requests that passed the proxy & Confirmed healthy under 100-VU sustained concurrent load & Application architecture is sound \\
    \hline
    Key implication & Rate-limiting thresholds must be re-tuned before go-live to reflect legitimate concurrent usage patterns & Baseline confirmed: application can sustain 100-VU load; tune rate-limiter, do not remove it & Production config action required \\
    \hline
  \end{tabular}
  }%
  \vspace{2ex}
  \captionsetup{type=table}
  \caption{Performance test key findings across Version 1.0 and Version 1.1}
  \label{tab:perf-findings}
\end{table}

\noindent Version~1.0 of the performance test revealed that the Nginx reverse proxy
rate limiting policy had been configured at a threshold below what
100~simultaneous virtual users required, causing the proxy to return
HTTP~429 (Too~Many~Requests) before requests ever reached the .NET~8
application server.  This was confirmed by disabling the rate limiting rules
for Version~1.1, after which the error rate dropped substantially.

This finding is an important operational significance: If the rate limiting layer is always removed, this means that we \testit{cannot} take an important security measure for granted: blocking denial of service and credential stuffing attacks. Instead, it's important to set the thresholds to allow for the anticipated max concurrent users and same time, block abnormal amounts of request rates. The performance testing exercise itself highlighted the benefits of performing this testing with realistic levels of concurrency, and uncovered a configuration issue that would have impacted user engagement during run time and would not have been found during functional or security testing.

The performance PDF appendices, Appendix Version 1.0 and Appendix Version 1.1, contain a copy of full numeric metrics (such as request counts, P50/P95 response time percentiles, throughput graphs and per second error-rate charts). 

% ----------------------------------------------------------
\subsection{User Acceptance Testing}
% ----------------------------------------------------------

User Acceptance Testing (UAT) was carried out for a specific task for each individual role-scoped user from a structured Google Form (\textit{Veracism UAT Form.pdf}).  Specific UAT items were assigned to a certain role, and blank columns are provided for roles that were not assigned. Unassigned columns are included only to allow for the scoring of columns that have been assigned and are not intentional. The UAT cycle included participation from four different roles: System Administrator, Tenant Administrator, IT Integrator / Developer, and External Verifier.

\subsubsection{UAT participants and scope}

Participants, roles and anticipated activities during UAT are summarized in Table~\ref{tab:uat-participants}.

\begin{table}[H]
  \vspace{2ex}
  \centering
  \footnotesize
  {%
  \setlength{\tabcolsep}{6pt}
  \renewcommand{\arraystretch}{1.3}
  \begin{tabular}{|>{\raggedright\arraybackslash}p{0.30\textwidth}|
                  >{\centering\arraybackslash}p{0.11\textwidth}|
                  >{\centering\arraybackslash}p{0.23\textwidth}|
                  >{\centering\arraybackslash}p{0.18\textwidth}|}
    \hline
    \textbf{Role} & \textbf{Participants} & \textbf{UAT Cases Assigned} & \textbf{Evaluations Recorded} \\
    \hline\hline
    System Administrator & 1 & UAT-01, UAT-02, UAT-03 & 3 \\
    \hline
    Tenant Administrator & 2 & UAT-04, UAT-05, UAT-06 & 6 \\
    \hline
    IT Integrator / Developer & 2 & UAT-07 & 2 \\
    \hline
    External Verifier & 5 & UAT-08 & 5 \\
    \hline
    \textbf{TOTAL} & \textbf{10} & \textbf{8 unique UAT IDs} & \textbf{16} \\
    \hline
  \end{tabular}
  }%
  \vspace{2ex}
  \captionsetup{type=table}
  \caption{UAT participant distribution by role and assigned task scope}
  \label{tab:uat-participants}
\end{table}

\subsubsection{UAT results by role}

Table~\ref{tab:uat-by-role} presents the pass, partial, and fail breakdown
per role group across all recorded evaluations.

\begin{table}[H]
  \vspace{2ex}
  \centering
  \footnotesize
  {%
  \setlength{\tabcolsep}{6pt}
  \renewcommand{\arraystretch}{1.3}
  \begin{tabular}{|>{\raggedright\arraybackslash}p{0.28\textwidth}|
                  >{\centering\arraybackslash}p{0.16\textwidth}|
                  >{\centering\arraybackslash}p{0.09\textwidth}|
                  >{\centering\arraybackslash}p{0.09\textwidth}|
                  >{\centering\arraybackslash}p{0.09\textwidth}|
                  >{\centering\arraybackslash}p{0.09\textwidth}|}
    \hline
    \textbf{Role} & \textbf{Evaluations} & \textbf{Passed} & \textbf{Partial} & \textbf{Failed} & \textbf{Pass \%} \\
    \hline\hline
    System Administrator & 3 & 2 & 1 & 0 & 66.7\% \\
    \hline
    Tenant Administrator & 6 & 6 & 0 & 0 & 100.0\% \\
    \hline
    IT Integrator / Developer & 2 & 2 & 0 & 0 & 100.0\% \\
    \hline
    External Verifier & 5 & 5 & 0 & 0 & 100.0\% \\
    \hline
    \textbf{OVERALL} & \textbf{16} & \textbf{15} & \textbf{1} & \textbf{0} & \textbf{93.8\%} \\
    \hline
  \end{tabular}
  }%
  \vspace{2ex}
  \captionsetup{type=table}
  \caption{UAT results by participant role (role-scoped; blank columns in unassigned roles excluded)}
  \label{tab:uat-by-role}
\end{table}

\noindent Three roles -- Tenant Administrator, IT Integrator, and External Verifier --
achieved a \textbf{100\% pass rate} across all their assigned evaluations.
The System Administrator role recorded the only partial result (UAT-01).
The participant confirmed that all dashboard summary statistics loaded
successfully without error, but noted the following usability suggestion:
\textit{``Summary of information that can be seen in the system were loaded
successfully without any error. Should include redirect links for every
summary (more details or view full report).''}  This feedback is captured as a low impact usability Reflections and Enhance and will be made available in the next development cycle when navigation links will be provided to users from each dashboard summary card that link to the detail or report view. Both User and configuration management (UAT-02 and UAT-03) passed completely.

\subsubsection{Key UAT findings}

The following behaviours were shared by the participant responses:

\begin{enumerate}
  \item \textbf{The services of verifying public documents are available to everyone.}
    All five External Verifier participants, UAT-08, indicated that document verification using CID is working without register, and one said it was “Excellent” in comments.

  \item \textbf{Easily accepted 2 Tenant upload and file management workflows.}
    The Tenant Administrator participants all reported passing UAT-04, UAT-05, and UAT-06. Tenants records, user assignments and API key configuration were properly updated in the UI and audit log.
    
  \item \textbf{IT Integrator API access is working.}
    The two IT Integrator participants confirmed that UAT-07 passed, concluding that the integration of the authenticated API upload and retrieval of metadata are as expected by programmatic consumers.

  \item \textbf{Administrator dashboard needs to be improved in usability.}
    The participant (the System Administrator) completed part of UAT-01. Every data presented showed no error and the participant gave the following constructive suggestion: “Should include redirect links for all summaries, should provide more details and/or present whole report.” This targeted feedback will be leveraged in the subsequent development iteration by providing navigation links from each dashboard summary card to the page or detail view of the specific card.

  \item \textbf{All 8 UAT IDs received at least one response.}
    All tasks defined across all roles has been addressed within participant cohort of 10, achieving full scenario coverage in the UAT cycle.
\end{enumerate}

% ----------------------------------------------------------
\subsection{Security Testing}
% ----------------------------------------------------------

There was a first pass for security evaluation followed by a security evaluation pass where targeted manual code review and penetration testing were conducted for every security flagged found in the initial security evaluation pass – The first pass was performed using Burp Suite Community Edition – DAST was used.

\subsubsection{Initial DAST scan overview}

The findings returned by Burp Suite scan for the staging endpoint were at various severities. Each category of finding and post-review disposition of the findings are summarized in Table~\ref{tab:dast-overview}.

\begin{table}[H]
  \vspace{2ex}
  \centering
  \footnotesize
  {%
  \setlength{\tabcolsep}{5pt}
  \renewcommand{\arraystretch}{1.3}
  \begin{tabular}{|>{\raggedright\arraybackslash}p{0.22\textwidth}|
                  >{\centering\arraybackslash}p{0.14\textwidth}|
                  >{\centering\arraybackslash}p{0.10\textwidth}|
                  >{\centering\arraybackslash}p{0.14\textwidth}|
                  >{\raggedright\arraybackslash}p{0.23\textwidth}|}
    \hline
    \textbf{Finding Type} & \textbf{DAST Sev.} & \textbf{Count} & \textbf{Confidence} & \textbf{Post-Review Disposition} \\
    \hline\hline
    SQL Injection & High & 4 & Tentative & \textbf{False positive}: EF Core ORM; proxy returned HTTP~403 for malformed paths \\
    \hline
    Reflected XSS & Medium & 1 & Tentative & \textbf{False positive}: Angular output encoding absorbed injected payloads \\
    \hline
    Missing security headers & Low & Several & Certain & Partial: HSTS, X-Frame-Options present; cache-control absent on some endpoints \\
    \hline
    Backup file disclosure & Informational & Multiple & Certain & No risk: all backup URL patterns 403-gated at IIS; no accessible files \\
    \hline
    Email address disclosure & Informational & Multiple & Certain & By design: user email in authenticated admin API responses \\
    \hline
    Cacheable HTTPS responses & Informational & Multiple & Certain & Low risk: affects read-only public endpoints only \\
    \hline
  \end{tabular}
  }%
  \vspace{2ex}
  \captionsetup{type=table}
  \caption{Burp DAST initial findings and post review disposition – 21 March 2026.}
  \label{tab:dast-overview}
\end{table}

\noindent The reflected XSS finding, as well as all four SQL Injection findings was judged to be false alerts on manual code review. No handwritten SQL anywhere in the back-end codebase, it relies on a single line of code query to Entity Framework Core object graph created using LINQ. However it was not the application running behind the scanner that mis-handled the SQL data on the input to it, but it was responses from the IIS/Nginx proxy to malformed URL paths that triggered the scanner's SQL-injection heuristics.

\subsubsection{Listed vulnerabilities and remediation}

\textbf{Five real vulnerabilities} prior to graded deployments were identified and confirmed by manual code review and target penetration testing, all Medium-High or lower. Before the deployment date all five of the following were remediated. The complete summary is included in the Table~\ref{tab:security-detail} which includes the remediation technique used to meet a finding.

\begin{table}[H]
  \vspace{2ex}
  \centering
  \footnotesize
  {%
  \setlength{\tabcolsep}{5pt}
  \renewcommand{\arraystretch}{1.3}
  \begin{tabular}{|>{\centering\arraybackslash}p{0.07\textwidth}|
                  >{\raggedright\arraybackslash}p{0.19\textwidth}|
                  >{\centering\arraybackslash}p{0.13\textwidth}|
                  >{\raggedright\arraybackslash}p{0.34\textwidth}|
                  >{\centering\arraybackslash}p{0.12\textwidth}|}
    \hline
    \textbf{ID} & \textbf{Vulnerability} & \textbf{Severity} & \textbf{Remediation Applied} & \textbf{Status} \\
    \hline\hline
    VUL-01 & Unbounded API Client Data Response & Low & Server-side pagination (\texttt{OFFSET/FETCH}); configurable page/pageSize; max cap of 100/page; response envelope updated & Remediated \\
    \hline
    VUL-02 & BOLA on API Client Records (cross tenant read) & Medium--High & Explicit tenant ownership check on GET/PUT/PATCH/DELETE; HTTP~403 for cross tenant reads; admin retains cross tenant access & Remediated \\
    \hline
    VUL-03 & Insufficient File Type Validation & Low--Medium & Allowlist of 18 permitted extensions via \texttt{FileValidationHelper}; applied before scanning pipeline; frontend \texttt{accept} attribute updated & Remediated \\
    \hline
    VUL-04 & BOLA on Report Records (cross tenant report access) & Medium--High & Tenant ownership check on \texttt{GET /reports/\{id\}}, \texttt{/latest}, and \texttt{/user/\{userId\}}; HTTP~403 for non-owner callers & Remediated \\
    \hline
    VUL-05 & Missing Tenant Isolation on File View/Download & Medium--High & \texttt{EnforceTenantOwnershipAsync(cid)} helper in \texttt{VerifyController}; HTTP~403 for non-owner; admin pass-through retained & Remediated \\
    \hline
  \end{tabular}
  }%
  \vspace{2ex}
  \captionsetup{type=table}
  \caption{Confirmed security vulnerabilities and remediation status}
  \label{tab:security-detail}
\end{table}

\noindent Three of the five confirmed vulnerabilities (VUL-02, VUL-04, VUL-05) were classified as ‘Broken Object-Level Authorisation (BOLA)'. The African American reason funds are protected behind BOLA are sensitive academic papers, transcripts, credentials and institutional identifiers, whose confidentiality is a major platform guarantee. At the controller level, there are tenant scoped ownership checks performed across all three BOLA paths, which have also closed.

\subsubsection{Security posture summary}

Table~\ref{tab:security-posture} presents a one-page assessment of the platform's security posture across key dimensions at the time of initial controlled deployment.

\begin{table}[H]
  \vspace{2ex}
  \centering
  \footnotesize
  {%
  \setlength{\tabcolsep}{5pt}
  \renewcommand{\arraystretch}{1.3}
  \begin{tabular}{|>{\raggedright\arraybackslash}p{0.31\textwidth}|
                  >{\centering\arraybackslash}p{0.12\textwidth}|
                  >{\raggedright\arraybackslash}p{0.42\textwidth}|}
    \hline
    \textbf{Security Dimension} & \textbf{Assessment} & \textbf{Notes} \\
    \hline\hline
    Critical / High vulnerabilities & \textbf{None} & No confirmed High or Critical findings at deployment \\
    \hline
    Medium--High vulnerabilities (BOLA) & \textbf{Remediated} & VUL-02, VUL-04, VUL-05: all three BOLA paths closed \\
    \hline
    SQL injection attack surface & \textbf{Eliminated} & EF Core ORM with parameterised queries throughout; no handwritten SQL \\
    \hline
    Authentication (Identity Provider) & \textbf{Strong} & Auth0 JWT Bearer tokens; \texttt{ValidateUserExistsFilter} syncs users into local DB \\
    \hline
    Authorisation model & \textbf{Adequate} & \texttt{AdminOnly}, \texttt{TenantOnly}, \texttt{AdminOrTenant} policies; inline controller checks \\
    \hline
    Transport security & \textbf{Strong} & HTTPS-only; HSTS headers enforced via Nginx \\
    \hline
    File upload safety & \textbf{Good} & ClamAV + YARA + PDF phishing analysis; file-type allowlist (18 extensions) \\
    \hline
    Audit trail & \textbf{Partial} & Admin audit export implemented; tenant scoped audit UI deferred (DEF-004) \\
    \hline
    Deferred findings & \textbf{2 items} & DEF-01 (partial BOLA on user enumeration) and DEF-02 (nullable warnings); neither is immediately exploitable \\
    \hline
  \end{tabular}
  }%
  \vspace{2ex}
  \captionsetup{type=table}
  \caption{There is a security posture assessment for the initial deployment, where this take place, last March of 2026.}
  \label{tab:security-posture}
\end{table}

% ----------------------------------------------------------
\subsection{Overall Discussion}
% ----------------------------------------------------------

Overall, the five streams of testing show a uniform picture of a system that meets its basic functional goals and infrastructure-level configuration issues that can be fixed prior to going to production scale-out locations.

The \textbf{100\% automated assertion pass rate} throughout the 182 assertions confirms that simply uploading your content, scanning for threats, then storing the content in a decentralised way and then later uploading a public CID based verification also works the way it should. The 30-case manual suite failed merely three times (with an overall 90.0\% full-pass rate), representing gaps in non-core subsystems, retry logic, audit tooling, and verified document verification logging, which do not impact the integrity or confidentiality of verified documents.

There was a real operational worry, the rate limiting configuration configuration that the performance test exercise brought to light, which would have had a negative impact on the user experience if it had not been identified during the performance test. The v1.1 test verified that the application layer can handle up to 100-VU concurrent load, and that there were no application-level faults present in v1.0; it was only due to wrong configuration of infrastructure. It's a situation that can be fixed by configuration changes.

After the initial controlled deployment of the three critical risks identified, documented, and remediated before the deployment, the security assessment validated that the three most significant risks were identified and documented, and were completely remediated. Embedding the DevSecOps workflow in the development process proved effective in discovering and fixing access-control issues with automated DAST scanning and manual code review. There are no reported Critical or High severity vulnerabilities at deployment, indicating that a general baseline level of security is suitable for controlled institutional use.

Here's another system-wide observation from both security and functional testing camps that tenant-ownership checks are being performed as a series of inline code in each controller today. This pattern is correct but can add to the possibility of omissions in subsequent endpoint additions. Putting the checks into a common base controller, or a dedicated ASP. The next development phase of .NET Core will include a recommended architectural upgrade to NetsLib AuthorisationHandler.

